\section{Corrections and Remaining Limitations}
\label{sec:issues}

The thesis workflow was checked against the raw calculation outputs and the notebooks. Table~\ref{tab:issues} lists what was found, what has been corrected in this version, and what remains open.

\begin{table}[t]
\caption{Problems found in the thesis workflow and their status in this version.}
\label{tab:issues}
\centering
\small
\begin{tabularx}{\textwidth}{lXX}
\toprule
\textbf{Area} & \textbf{Problem} & \textbf{Status} \\
\midrule
Descriptors & Parser read the wrong lines of the ORCA output (Section~\ref{sec:issue-desc}). & Corrected; both descriptor sets compared in Table~\ref{tab:grouped}. \\
Structures & Six SMILES described a neutral molecule or an ion with its counter-ion, so ORCA treated them as doublet radicals. & Corrected and recalculated. \\
Convergence & [PF$_6$]$^-$ dissociated into PF$_5$ + F$^-$ during its optimisation; [HSO$_4$]$^-$ and [C$_3$F$_7$CO$_2$]$^-$ had one imaginary frequency. & Recalculated; all training ions are minima (Section~\ref{sec:comp}). \\
Validation & Random 80/20 split of points; most test ILs also appear in training. & Replaced by grouped cross-validation (Section~\ref{sec:grouped}). \\
Level of theory & One stray input file (\path{calculations/ions/BMIM.inp}) states a B3LYP/def2-SVP single point. & Not used; all descriptors are r2SCAN-3c. \\
Descriptor physics & Gas-phase isolated ions, one conformer, no diffuse functions. & Open: corrected descriptors match atom counts but do not beat them (Section~\ref{sec:grouped}). \\
Screening & Done with the thesis model. & Open: to be recomputed. \\
Bio-based ILs & Laurate computed with B97-3c; typed-in values differ from outputs; atom counts entered under wrong column names. & Open: predictions not reliable (Section~\ref{sec:issue-bio}). \\
Data source & Source of \path{data/raw/rawdata.xlsx} not cited; no reference per point. & Identified (Laakso et al.\cite{laakso2025machine}, data of Song et al.\cite{song2020prediction}); 78\% of points traced to a publication; main results use the referenced dataset. \\
Duplicates & Table S5 is a subset of Table S6, so the thesis set repeats points: 2,013 of 11,256 rows, and 650 of the 2,252 random-split test points (29\%) have an identical copy in training. & Open for the thesis protocol (Table~\ref{tab:metrics} overstates accuracy further); the referenced dataset has no repeated point from the same publication. \\
References & DOIs of the referenced dataset taken from an unchecked search: 51 of 171 publications (4,728 points) had a wrong or missing DOI. & Corrected and verified against Crossref (Section~\ref{sec:data}). \\
\bottomrule
\end{tabularx}
\end{table}

\subsection{Descriptor extraction in the thesis models}
\label{sec:issue-desc}
The parser in \path{notebooks/thesis/07-descriptor-extraction.ipynb} was compared with the ORCA outputs for the 76 ions in the thesis set. Three errors were found:
\begin{itemize}
    \item \textbf{Total energy} was the first \texttt{FINAL SINGLE POINT ENERGY} in the file, which belongs to the starting geometry, not the optimised one. The median difference from the optimised energy is 0.004~Hartree, but it reached 0.92~Hartree for [PF$_6$]$^-$.
    \item \textbf{HOMO energy} was the largest number found in any four-column table after the last orbital block, not only in the orbital block itself. For 39 of the 76 ions the stored HOMO was a value from an unrelated table and lay above the LUMO; for example, [BF$_4$]$^-$ had a stored HOMO of $+0.657$~Hartree against $-0.110$~Hartree in the output. The LUMO was correct for 74 ions.
    \item \textbf{Dipole moment} was the $x$ component of the dipole vector in atomic units, not the magnitude in Debye. When that component was negative the pattern did not match, the value was left empty and later filled with zero; this happened for 44 of the 76 ions.
\end{itemize}
The corrected parser (\path{scripts/extract_descriptors.py}) reads the last energy, the last orbital-energy block and the dipole magnitude in Debye. Its HOMO energies follow the published ordering of ion ionisation energies, whereas the thesis values do not (Section~\ref{sec:desc-check}).

\subsection{Limits of isolated-ion descriptors}
The corrected descriptors are computed for one conformer of each ion in the gas phase. They do not describe the interaction between cation and anion, the interaction of either ion with \co{}, or the structure of the liquid, all of which control physical solubility. The r2SCAN-3c basis set has no diffuse functions, which makes gas-phase anion orbital energies less reliable. Consistent with this, the corrected descriptors predict unseen anions as well as elemental atom counts do, but not better (Table~\ref{tab:grouped}).

\subsection{Bio-based IL inputs}
\label{sec:issue-bio}
The oleate and laurate predictions in Section~\ref{sec:screening} have four problems. First, the oleate anion was calculated with r2SCAN-3c like the training ions, but the laurate anion was calculated with B97-3c\cite{brandenburg2018b973c}. Second, the descriptor values typed into the prediction notebook do not all match the output files. The laurate total energy entered is $-1516.05$~Hartree against $-621.36$~Hartree in the output, and the dipole moments entered (13.5~D and 7.3~D) differ from the output magnitudes (34.9~D and 18.6~D). Third, the atom counts were entered under the names \texttt{anion\_C}, \texttt{anion\_H} and \texttt{anion\_O}, which do not exist in the model. The missing columns, including all cation counts, were therefore filled with zero. Fourth, imidazolium carboxylates can react with \co{} through the carbene formed at the C2 position of the ring, which a physical-solubility model does not describe.
